%% chapters/behavior_modeling/ch2_correlation.tex
%% 第2章　相関・異質性
\section{MNLからの拡張: 相関・異質性}\label{sec:abms}
本節では，前節で導入したMNLモデルからの拡張として，選択肢間の相関や意思決定主体間の異質性を考慮したモデルを解説する．
具体的には，選択肢間の相関を部分的に許すNested Logitモデル，個人の異質性を考慮するMixed Logitモデル，そして経路選択特有の拡張であるPath-Size Logitモデルなどを紹介する．

\subsection{選択肢間の相関とIIAの限界}
まず選択肢間の相関の問題を考える．


% 第3節 MNLからの拡張：相関・異質性

% ここで学生が最初につまずくので、拡張理由を明確に書くべきです。

% 3.1 IIAとその限界
% 類似選択肢があると不自然
% 経路の重複問題
% 3.2 Nested Logit
% 誤差相関を部分的に許す
% 木構造による整理
% 3.3 Mixed Logit
% 係数異質性
% 柔軟性
% シミュレーション推定
% 3.4 Path-Size Logit など経路選択特有の拡張
% 経路重複への対応
% 経路選択の相関は交通分野で特有に重要

% 3.5
% 観測異質性と非観測異質性
% 交互作用項で表せること／表せないこと

% ここで大事なのは、
% Nested Logit は「選択肢相関」への対応、Mixed Logit は「個人異質性 + 柔軟な相関」への対応
% と役割を分けて説明することです
